\subsection{Stopped depth for a general finite complex network}
\label{sec:general-network-stopped-depth}
This lemma is conditional on a finite network's exact-arithmetic recurrence;
it does not supply a new network or multiplication exponent.
Let $m\ge3$, $s\ge2$, and let $E\ge0$ be a verified additive elementary-operation
charge at each internal node. Suppose
\[
 A(e)\le sA(e/m)+E,\qquad A(e)\le8e\quad(e<d^\beta).
\]
Choose rational $\alpha\ge1$, $0<\beta<1$, $\zeta>0$ with $s<m^\alpha$.
For $\alpha=u/v$, this hypothesis is the exact integer comparison $s^v<m^u$.
Set
\[
 B=s+E,\quad C_1=\alpha-(\alpha-1)\beta+\zeta,
\quad C_0=\left\lceil\max\{128mB^2,18mB^2(1+1/\zeta)\}\right\rceil.
\]
For a root $e\le d$, the number $j$ of internal levels is at most
$(1-\beta)\log_m d+1$. Thus
\[
 s^j\le s d^{\alpha(1-\beta)},\qquad
 A(e)\le8d^\beta s^j+E\sum_{i<j}s^i
 \le s(8+E)d^{\alpha-(\alpha-1)\beta}
 \le9B^2d^{\alpha-(\alpha-1)\beta}.
\]
The same bound covers a root that is already a leaf. Since $s\ge2$, the
geometric overhead sum is at most $s^j$; $E$ remains additive rather than
being inserted into the branching factor. Also $\alpha\ge1$ makes the last
exponent at least one.

There are at most $m(1+1/\zeta)d^\zeta$ base-$m$ pieces. Retaining the
established $18d$ allowance for individual axes and outer phases gives
\[
 A_{\rm layer}\le9mB^2(1+1/\zeta)d^{C_1}+18d
 \le C_0d^{C_1}.
\]
Consequently the existing denominator and magnitude induction works with
$\Delta=\lceil C_0d^{C_1}\rceil$, and the coefficient width remains $O(p)$
when $\epsilon C_1<1$. This replaces the special hypothesis $s<m^5$;
taking $\alpha=5$ recovers the retained expression $C_1=5-4\beta+\zeta$.
New scalar coefficients must be compiled into the counted elementary
operations and included in $E$. Neither the phase-kernel hypotheses nor the
finite network's scalar correctness follow from this depth lemma.

\paragraph{Interaction with the finite savings.}
Put $t=1-\beta$. The retained Gaussian and layer constraints imply
\[
 \kappa<\epsilon\min(a_{\rm b},ta_{\rm c}),\qquad
 \epsilon<\min\{1/5,1/[1+(\alpha-1)t]\}.
\]
Therefore this subset of necessary conditions gives the upper bound
\[
 \kappa\le
 \frac{\min(a_{\rm b},a_{\rm c})}
 {\max\{5,1+(\alpha-1)\min(1,a_{\rm b}/a_{\rm c})\}}.
\]
To see this, for $t\le a_{\rm b}/a_{\rm c}$ the quotient
$ta_{\rm c}/\max(5,1+(\alpha-1)t)$ is increasing. Above that point its
numerator is constant and its denominator is nondecreasing. The maximizing
limit is $t=\min(1,a_{\rm b}/a_{\rm c})$. This is an upper bound from selected
constraints, not an assertion that all assembly requirements attain it.
In particular, this bound retains the value $a_{\rm b}/5$ if
\[
 a_{\rm c}/a_{\rm b}\ge\max\{1,(\alpha-1)/4\}.
\]
Larger role counts can thus require more complex-network headroom without
forcing a quadratic dependence on the finite saving.
